\documentclass[11pt]{article}
\usepackage[margin=1in]{geometry}
\usepackage{amsmath,amssymb}
\usepackage{enumitem}
\usepackage[hidelinks]{hyperref}
\usepackage[T1]{fontenc}
\usepackage[utf8]{inputenc}
\usepackage{booktabs}
\usepackage{tabularx}
\usepackage{array}

\title{Precision Routing for Efficient Robot Policy Deployment on Edge Devices\\
\large CS 2420 Project Plan}
\author{Brianna Grissom, Hang Le, Shresth Verma, Lester Xue}

\begin{document}
\maketitle

\section*{Research assignments}

\noindent
Brianna:

\noindent
Lester:

\noindent
Shresth:

\noindent
Hang:




\section*{Classification model to switch precision per step}\label{classification-model-to-switch-precision-per-step}

We want to build a classification model that determines whether to switch precision at each action step during deployment.

\subsection{Classifier features}\label{classifier-features}

\begin{itemize}
\item The classifier itself has tiny compututation.The computationally expensive part is obtaining the features. If a feature requires high computation, it costs more than the savings from low precision.
\item So we prefer pre-computed features or features that come from previous steps and ideally overlapped with the GPU's image-processing phase (as DyQ-VLA does on the CPU).
\end{itemize}

\subsubsection{Candidate features (table 1)}\label{candidate-features}

\begin{table}[h]
\centering
\small
\begin{tabularx}{\textwidth}{>{\raggedright\arraybackslash}X >{\raggedright\arraybackslash}p{0.2\textwidth} >{\raggedright\arraybackslash}p{0.12\textwidth} >{\raggedright\arraybackslash}X}
\toprule
\textbf{Feature} & \textbf{Source} & \textbf{Cost} & \textbf{Notes} \\
\midrule
Previous action(s): translation xyz, rotation, gripper command & Model's own past outputs & Free & Short history window (e.g.~last 5--10 steps) \\
Action changes between steps (translational and rotational change) & Computed from action history & Negligible & Captures sudden twists and slowdowns \\
Motion Fineness and Angular Jerk (from DyQ-VLA paper) & Computed from action history & Negligible & Useful baseline features; lets us compare against the heuristic directly \\
Joint positions and velocities & Robot proprioception & Free & Already read every control step \\
End-effector pose and speed & Proprioception / forward kinematics & Negligible &  \\
Gripper state (open/closed, width) & Proprioception & Free & Strong cue for grasp and release phases \\
Token confidence: entropy and top-1 vs top-2 margin per action token & Previous step's output distribution & Free & OpenVLA-style discrete action tokens give these for free \\
Episode progress (step t / expected max steps) & Step counter & Free & ElegantVLA uses this as a coarse temporal cue \\
Previous precision decisions & Dispatcher state & Free & Helps the classifier learn smooth switching \\
Force/torque or contact signals & Wrist sensor, if available & Free (if hardware has it) & Possibly the most direct contact cue; not available on every robot \\
Task instruction embedding & Computed once per episode & Negligible per step & Lets the classifier learn task-specific behaviour \\
\bottomrule
\end{tabularx}
\caption{Candidate classifier features.}
\end{table}

\subsubsection{Expensive features}\label{expensive-features}

\begin{itemize}
\item \textbf{Separate vision encoder or extra image processing}: adds a second expensive pass.
\item \textbf{Hidden-state similarity across steps} (as in ElegantVLA's first-layer check): this is only cheap if the states are already cached; check actual latency.
\item \textbf{Low-precision proposed action for the current step} (cascade design): fine as a feature, but every rejection costs a second pass. Only pays off if rejections stay below roughly 1 $-$ (low-precision cost / full-precision cost), i.e.~about 33\% at a 1.5$\times$ speedup.
\end{itemize}

\subsection{What to quantize: activations, weights, or both}\label{what-to-quantize-activations-weights-or-both}

A model has two kinds of numbers we could store with less bits:

\begin{itemize}
\item \textbf{Weights}: what the model learned in training. They are stored permanently and read again at every step.
\item \textbf{Activations}: the working numbers the model computes from the current camera image and instruction. They are created fresh at each step and then thrown away.
\end{itemize}

An OpenVLA step has two phases, and each is slowed down by something different:

\begin{itemize}
\item \textbf{Image + instruction processing}: hundreds of tokens are handled at once. \textbf{The arithmetic is the slow part.}
\item \textbf{Action generation}: 7 action tokens are generated one at a time, and each one reads every weight. \textbf{Reading weights from memory is the slow part}. Rough estimate on an A100: $\sim$7--9 ms per token at FP16, so $\sim$50--60 ms for all 7. At 4-bit weights it's about a quarter of that.
\end{itemize}

This matters because lowering weight precision mostly speeds up action generation, while lowering activation precision mostly speeds up image processing.

\subsubsection{Option 1: Switch activation precision (DyQ-VLA's approach)}\label{option-1-switch-activation-precision-dyq-vlas-approach}

\textbf{Weights stay fixed} (DyQ-VLA keeps them at 4 bits), and each step picks 2, 4 or 8 bits, or full precision, for the activations.

\textbf{Benefits}

\begin{itemize}
\item \textbf{Switching precision costs nothing}: activations are recomputed every step anyway, so there's nothing extra to store or load.
\item It has already been shown to work. DyQ-VLA gets $\sim$1.5$\times$ speedup while keeping 99.5\% of performance.
\item It speeds up the image phase, if the hardware has fast low-precision math.
\end{itemize}

\textbf{Problems}

\begin{itemize}
\item \textbf{Almost no memory saving, since activations are temporary}. DyQ-VLA's memory savings come from its fixed 4-bit weights, not from the switching.
\item \textbf{Limited speedup in action generation, which is slowed by reading weights rather than by arithmetic}. This is likely why the speedup tops out around 1.5$\times$.
\item It needs custom low-precision math code. DyQ-VLA's switching first added 16.8 ms per step until they wrote custom kernels and moved the decision logic onto the CPU.
\item Activations can be hard to round accurately. In language models they often contain a few extreme values that make them harder to round than weights.
\end{itemize}

\subsubsection{Option 2: Switch weight precision (nested weights)}\label{option-2-switch-weight-precision-nested-weights}

Store the weights once at the highest precision we want (e.g., 8 bits) in a ``nested'' layout, where the lower-precision versions are contained inside the higher-precision one. Each weight's first few bits give a rough value, and more bits refine it, like reading more digits of a price. The bits are stored grouped by position (the first bit of every weight together, then the second bit of every weight, and so on), so running at lower precision just means reading fewer groups. Prior work: \href{https://arxiv.org/pdf/2402.10517}{Any-Precision LLM} (ICML 2024) and \href{https://arxiv.org/abs/2502.06786v2}{Matryoshka Quantization} (ICML 2025).

\textbf{Related work} (details in ``Related work'' below)

\begin{itemize}
\item \textbf{Fixed 4-bit weights already work for OpenVLA.} QAIL and SQIL fine-tune OpenVLA to 4-bit weights and lose about 1 point of success, while running 2.5$\times$ faster on a Jetson AGX Orin (15.2 $\rightarrow$ 4.0 GB). QVLA and Mix-QVLA show that 4-bit weights with full-precision activations match the full-precision success rate without fine-tuning (76.5\% and 76.6\% vs.~76.5\%). So an 8-bit top level is probably unnecessary. \textbf{We could store 4 bits as the top level and drop to 2--3 bits on easy steps, using about the same memory as DyQ-VLA.}
\item \textbf{MotionVQ} is the closest prior work. It switches between a coarse and a fine set of weights every step based on how much the arm moved. But it stores two separate sets rather than nested ones, uses a fixed motion threshold, and reports speed only on a simulated custom chip. According to the summary, it uses the coarse weights during \emph{low}-motion steps, the opposite of DyQ-VLA's assumption, so check the paper.
\item \textbf{MoBiQuant} (language models, not robotics) uses nested weights and picks a precision for each token with a router. It's the closest example of nested weights with runtime switching.
\item \textbf{BitVLA} shows weights can go as low as $-$1/0/1 if the model is trained that way from scratch, so very low weight precision isn't hopeless for VLAs.
\end{itemize}

\textbf{Benefits}

\begin{itemize}
\item It targets the main bottleneck in action generation, so \textbf{it could beat Option 1's speedup.}
\item \textbf{Only one copy of the model is stored}, and switching just means reading more or fewer bit groups.
\item \textbf{It makes a cascade much cheaper}. Rerunning only the action tokens at higher precision costs just the extra bit groups, not a second full pass, so the $\sim$33\% break-even estimate above should be redone.
\end{itemize}

\textbf{Problems}

\begin{itemize}
\item \textbf{The highest available precision is whatever we store}, so there's no FP16 unless we pay the memory for it. We need to check that 8-bit weights alone keep OpenVLA's success rate.
\item It uses more memory than DyQ-VLA: $\sim$7 GB at 8 bits vs.~their 4.7 GB, though still far below FP16's 15.2 GB.
\item Simply chopping off bits could give a worse low-bit model than one purposely made for that bit-width. This hasn't been tested on VLAs.
\item It needs special code that reads the bit groups. Any-Precision LLM released code for LLaMA-style models (OpenVLA's backbone is Llama 2), but it may need adapting.
\item \textbf{It does little for the image phase}, which is slowed by arithmetic.
\item Very low weight precision (2-bit) may lose a lot of accuracy and may only be usable during the easiest parts of a task.
\end{itemize}

Simpler alternatives that don't work:

\begin{itemize}
\item \textbf{Keep a separate copy at each precision on the GPU.} 2 + 4 + 8 + 16 bits is $\sim$30 bits per weight, which is more memory than FP16 alone.
\item \textbf{Load copies from the CPU when needed.} Moving a 4-bit 7B model ($\sim$3.5 GB) over PCIe takes $\sim$140 ms, longer than a whole control step.
\end{itemize}

\subsubsection{Option 3: Switch both}\label{option-3-switch-both}

\textbf{Use weight precision for the action-generation phase and activation precision for the image phase, so each targets the phase where it's the bottleneck.}

\textbf{Related work} (details in ``Related work'' below)

\begin{itemize}
\item \textbf{Fixed weights + activations at 4 bits already keeps around the same success rate.} QVLA (OpenVLA: 76.0\% vs.~76.5\%, 4.3 GB, 1.47$\times$), Mix-QVLA (OpenVLA: 76.3\% vs.~76.5\%, 4.0 GB, 1.52$\times$), QuantVLA, HoloQ-VLA and FoldQuantVLA all keep success within about a point of full precision at 4-bit weights and 4- or 8-bit activations. \textbf{This is the bar Option 3 has to beat}. Matching static 4-bit isn't enough: \textbf{the gain must come from going \emph{below} 4 bits during easy steps.}
\item \textbf{No paper switches both per step.} DyQ-VLA switches only activations and MotionVQ only weights.
\item \textbf{FoldQuantVLA} found that keeping a few sensitive layers at 8 bits raised real-robot success from 80.0\% to 92.5\%. \textbf{Fixed per-layer exceptions like this could be combined with per-step switching for the remaining layers.}
\item \textbf{BitVLA} ($-$1/0/1 weights, 8-bit activations, 1.4 GB, 96.0\% on LIBERO) shows how far fixed precision for both can go when the model is trained for it.
\end{itemize}

\textbf{Benefits}

\begin{itemize}
\item \textbf{The largest potential speedup, since both phases get faster.}
\item \textbf{The two can be controlled separately.} For example, the image phase could always run with low-precision activations while the classifier only chooses weight precision for the action tokens, which keeps the decision simple.
\end{itemize}

\textbf{Problems}

\begin{itemize}
\item \textbf{Errors from both sources add up}, and the combined effect on action quality is untested.
\item \textbf{The classifier has more choices} (pairs of settings), which means more labels and rollouts and harder calibration. We could limit it to a few preset combinations.
\item \textbf{It needs the most engineering}: two sets of low-precision code and two sources of switching overhead.
\end{itemize}

\subsubsection{Open questions to test}\label{open-questions-to-test}

\begin{itemize}
\item Does OpenVLA keep its success rate with 8-bit weights and full-precision activations? This sets the highest precision for Option 2.
\item At each bit-width, \textbf{which hurts OpenVLA's action accuracy more: rounding weights or rounding activations?}
\item \textbf{On our hardware, how is the wall-clock time per step split between image processing and action generation? This decides which option has more to gain.}
\item How much does the image phase's precision affect action quality?
\end{itemize}

\subsection{Related work}\label{related-work}

\subsubsection{Fixed precision, weights only}\label{fixed-precision-weights-only}

\begin{itemize}
\item \textbf{\href{https://arxiv.org/pdf/2412.01034}{QAIL}} (Quantization-Aware Imitation Learning): fine-tunes OpenVLA to 4-bit weights. 73.1\% vs.~74.0\% on LIBERO, with 2.5$\times$ speedup and energy savings on a Jetson AGX Orin.
\item \textbf{\href{https://arxiv.org/abs/2505.15304}{SQIL}} (Saliency-Aware Quantized Imitation Learning): same setting as QAIL (OpenVLA and $\pi$0, 4-bit weights). 73.2\% vs.~73.8\% on LIBERO, 2.5$\times$ faster on Jetson AGX Orin, 15.2 $\rightarrow$ 4.0 GB. During fine-tuning it weights the robot states where small action errors matter most (top 20\% by perturbation sensitivity) more heavily. That's a fixed-precision cousin of our per-step sensitivity labels.
\end{itemize}

\subsubsection{Fixed precision, weights and activations}\label{fixed-precision-weights-and-activations}

\begin{itemize}
\item \textbf{\href{https://arxiv.org/abs/2602.03782}{QVLA}} (ICLR 2026): each output channel gets its own bit-width from \{0, 2, 4, 8, 16\}, set once offline. OpenVLA-OFT at W4A4: 96.0\% vs.~97.1\%, 4.5 GB, 1.49$\times$. OpenVLA at W4A4: 76.0\% vs.~76.5\%, 4.3 GB, 1.47$\times$. OpenVLA at W4A16 matches full precision. This is the main baseline DyQ-VLA compares against.
\item \textbf{\href{https://arxiv.org/abs/2606.19565}{Mix-QVLA}}: one bit-width per layer from \{2, 4, 8, 16\}, chosen offline by optimization. Fixed at deployment: it ``does not perform timestep-wise bit switching during deployment.'' OpenVLA at W4A4: 76.3\% vs.~76.5\%, 4.0 GB, 1.52$\times$. W4A16: 76.6\%.
\item \textbf{\href{https://arxiv.org/abs/2602.20309}{QuantVLA}}: W4A8 on $\pi$0.5 and GR00T N1.5, no training. $\pi$0.5: 97.6\% vs.~97.1\%, 4.27 $\rightarrow$ 1.28 GB.
\item \textbf{\href{https://arxiv.org/abs/2605.28803}{HoloQ-VLA}}: uniform W4A4 on $\pi$0.5 and GR00T N1.5, including the diffusion action head. $\pi$0.5: 98.0\% vs.~97.1\%.
\item \textbf{\href{https://arxiv.org/abs/2609.24433v1}{FoldQuantVLA}}: W4A4 on GR00T and $\pi$0.5, with 1.2--1.5$\times$ speedups on real hardware, including a Jetson AGX Orin. Keeping a few sensitive layers at W8A8 raised real-robot success from 80.0\% to 92.5\%.
\item \textbf{\href{https://arxiv.org/abs/2506.07530}{BitVLA}}: a new 3B VLA trained with $-$1/0/1 weights and 8-bit activations. 1.4 GB, 96.0\% on LIBERO. It's trained from scratch rather than being a compressed existing model.
\end{itemize}

\subsubsection{Per-step (timestep-wise) switching}\label{per-step-timestep-wise-switching}

\begin{itemize}
\item \textbf{\href{https://arxiv.org/pdf/2603.07904}{DyQ-VLA}}: weights fixed at 4 bits. Activation precision switches every step (2/4/8 bits or full precision) based on a motion score.
\item \textbf{\href{https://pith.science/paper/2607.24148}{MotionVQ}} (arXiv 2607.24148, July 2026): weights only, using vector quantization. It keeps a coarse and a fine set of weights and switches between them every step based on motion in the previous step. Tested on five VLAs (OpenVLA, OpenVLA-OFT, RDT, $\pi$0, GR00T), with $\sim$2.5\% success loss. Speedups come from a simulated custom chip, not a real GPU. I've only read a summary, not the paper.
\end{itemize}

\subsubsection{Outside robotics: nested weights}\label{outside-robotics-nested-weights}

\begin{itemize}
\item \textbf{\href{https://arxiv.org/pdf/2402.10517}{Any-Precision LLM}} (ICML 2024) and \textbf{\href{https://arxiv.org/abs/2502.06786v2}{Matryoshka Quantization}} (ICML 2025): store one copy of the weights that can run at several precisions.
\item \textbf{\href{https://arxiv.org/abs/2602.20191}{MoBiQuant}} (2026): nested weights, with a router that picks the precision for each token at inference. Up to 1.34$\times$ faster than earlier any-precision methods. The paper page doesn't say how the router is trained.
\end{itemize}

\subsubsection{Who decides when to switch?}\label{who-decides-when-to-switch}

\textbf{None of the robotics papers above uses a supervised (trained) model to decide when to switch precision.}

\begin{itemize}
\item DyQ-VLA uses a hand-designed score (Motion Fineness + Angular Jerk) with fixed cutoffs.
\item MotionVQ calls its switch a ``state predictor,'' but the summary I read describes a fixed motion threshold (0.8). Check the paper to confirm nothing in it is learned.
\item MoBiQuant has a per-token router, but in language models, and its training wasn't checked.
\end{itemize}

\textbf{Gap for our project:} a learned (supervised) classifier that decides per step, possibly combined with nested weights. This field moves fast (several of these papers are from the last few months), so search again before committing.

\end{document}
